Observability for federated and distributed ML sits at the intersection of two bodies of work: general-purpose experiment tracking systems and the monitoring facilities built into individual FL frameworks. This section reviews both, and positions HiveWatch as a framework-agnostic layer that sits between them.

\textbf{General-purpose experiment tracking.} Tools such as Weights \& Biases~\cite{wandb}, MLflow~\cite{mlflow}, and TensorBoard~\cite{tensorboard} are widely used to log scalar metrics, hyperparameters, and artifacts for centralized training. These tools are effective at visualizing a single model's training curve, but they have no native notion of a federated round, a client population, per-client staleness, or client geography. Adapting them to FL typically means hand-rolling a schema for per-client metrics (e.g., manually namespacing keys such as \texttt{client/\textless id\textgreater/accuracy}), recomputing communication and divergence statistics in every project, and choosing a single backend up front — since these tools are not designed to fan the same event out to multiple destinations simultaneously. HiveWatch is built directly on top of this gap: it defines the FL-specific schema once (Section~\ref{sec:algorithm}) and uses Weights \& Biases and MLflow purely as emitter backends, so existing dashboards and team workflows built around either tool continue to work without modification.

\textbf{Monitoring inside FL frameworks.} Modern FL frameworks ship their own logging facilities, but these are tightly coupled to the framework's runtime and communication stack. APPFL~\cite{ryu2022appfl, li2025advances} and FedML~\cite{he2020fedml} expose configuration-driven logging that is primarily oriented around server-side console and file logs. Flower~\cite{beutel2020flower} provides a \texttt{Metrics} aggregation hook and an optional telemetry integration, but per-client visualization and replay are left to the application. NVIDIA FLARE~\cite{roth2022nvidia} integrates with TensorBoard and MLflow through framework-specific receivers tied to its job and controller abstractions. Flotilla~\cite{banerjee2025flotilla} targets heterogeneous resource scheduling rather than observability. In each case, switching FL frameworks means re-learning a different logging surface and, in practice, re-implementing custom dashboards for per-client diagnostics such as straggler detection, communication cost breakdowns, and client geography, none of which are portable across frameworks. HiveWatch instead exposes the same three-call runtime API (Section~\ref{sec:algorithm}) regardless of the underlying framework, and we demonstrate this directly by integrating the identical API into APPFL, Flower, and NVIDIA FLARE with no changes to HiveWatch itself (Section~\ref{sec:eval}).

\textbf{Heterogeneity, stragglers, and asynchronous FL.} A large body of FL research addresses client heterogeneity and stragglers at the \emph{algorithmic} level, through selective client participation~\cite{lai2021oort, wolfrath2022haccs}, adaptive local computation~\cite{li2020federated, diao2020heterofl, li2024fedcompass}, layer-wise aggregation~\cite{lang2024stragglers}, or asynchronous protocols that tolerate stale updates~\cite{xie2019asynchronous, nguyen2022federated, iakovidou2024asynchronous}. These methods change how the FL algorithm behaves in the presence of heterogeneity, but none of them are designed to help a practitioner \emph{observe} that heterogeneity in the first place — to see which specific clients are slow, by how much, and whether that correlates with geography, payload size, or system load. HiveWatch is complementary to this line of work: it records the staleness, gradient divergence, and per-client timing signals that motivate algorithms such as FedProx~\cite{li2020federated} and FedCompass~\cite{li2024fedcompass} in the first place, independent of which aggregation strategy a given run uses.
